\section{Notation}
\label{appx_sec:notation}

We summarize the main mathematical symbols and their meanings in Table~\ref{appx_tab:notation}.
Symbols are organized into six functional categories (highlighted in grey) following the two-layer structure of \projectname.
The first three cover the \textbf{RSPL substrate}: indexing conventions, the resource entity tuple, and protocol-registered resources including the context manager $\mathcal{M}_\tau$ and server interface $\mathcal{A}_\tau$.
The remaining three cover the \textbf{SEPL layer}: evolvable variables and the trainable subspace $\Theta$, the auxiliary spaces ($\mathcal{P}$, $\mathcal{Z}$, $\mathcal{H}$, $\mathcal{D}$, $\mathcal{G}$, $\mathcal{S}$) and five canonical reflection operators $\{\rho, \sigma, \iota, \varepsilon, \kappa\}$, and iteration-level variables of the evolutionary loop.

\begin{table}[H]
  \centering
  \caption{Notation used in the paper. Grey rows indicate categories.}
  \label{appx_tab:notation}
  \footnotesize
  \renewcommand{\arraystretch}{0.85}
  \setlength{\tabcolsep}{3pt}
  \begin{tabular}{p{3.2cm}p{9.2cm}}
    \toprule
    \textbf{Symbol} & \textbf{Description} \\
    \midrule
    \rowcolor{gray!15}
    \multicolumn{2}{c}{\textit{\textbf{Indexing and Sets}}} \\
    $\Tau$ & Set of RSPL entity types, $\{\textsc{Prompt},\textsc{Agent},\textsc{Tool},\textsc{Env},\textsc{Mem}\}$. \\
    $\tau$ & Entity type index, $\tau \in \Tau$. \\
    $\mathcal{I}_{\tau}$ & Index set of resource instances of type $\tau$. \\
    $i$ & Instance index, $i \in \mathcal{I}_{\tau}$. \\
    $\mathbb{V}$ & Space of version strings. \\
    $\wp(\cdot)$ & Power set operator. \\
    \midrule
    \rowcolor{gray!15}
    \multicolumn{2}{c}{\textit{\textbf{RSPL Resource Entity (Def.~\ref{appx_def:resource-entity})}}} \\
    $e_{\tau,i}$ & Resource entity tuple $(n_{\tau,i}, d_{\tau,i}, \phi_{\tau,i}, g_{\tau,i}, m_{\tau,i})$. \\
    $n_{\tau,i}$ & Unique resource name. \\
    $d_{\tau,i}$ & Short description. \\
    $\phi_{\tau,i}:\mathcal{X}_{\tau}\!\rightarrow\!\mathcal{Y}_{\tau}$ & Input-to-output mapping of the resource. \\
    $g_{\tau,i}$ & Evolvability marker, $g_{\tau,i} \in \{0,1\}$, indicating whether the resource is evolvable. \\
    $m_{\tau,i}$ & Auxiliary metadata dictionary. \\
    $\mathcal{E}_{\tau}$ & Set of resource entities of type $\tau$. \\
    \midrule
    \rowcolor{gray!15}
    \multicolumn{2}{c}{\textit{\textbf{RSPL Registration Record (Def.~\ref{appx_def:resource-registration-record})}}} \\
    $c_{\tau,i}$ & Registration record $(e_{\tau,i}, v_{\tau,i}, \eta_{\tau,i}, \theta_{\tau,i}, \mathcal{F}_{\tau,i})$. \\
    $\mathcal{C}_{\tau}$ & Set of registration records for type $\tau$. \\
    $v_{\tau,i}$ & Version string of the resource instance. \\
    $\eta_{\tau,i}$ & Implementation descriptor (e.g., import path, class, or source). \\
    $\theta_{\tau,i}$ & Instantiation parameters (e.g., constructor arguments). \\
    $\mathcal{F}_{\tau,i}$ & Exported representations for LLM interaction (schemas/text/structured args). \\
    \midrule
    \rowcolor{gray!15}
    \multicolumn{2}{c}{\textit{\textbf{Protocol-registered Resource (Def.~\ref{appx_def:protocol-registered-resource})}}} \\
    $\mathcal{R}_{\tau}$ & Type-specific registry of protocol-registered resources. \\
    $\mathcal{R}$ & Global registry, $\bigcup_{\tau}\mathcal{R}_{\tau}$. \\
    $\mathcal{M}_{\tau}$ & Context manager for type $\tau$ (maintains registry and version lineage). \\
    $\mathcal{A}_{\tau}$ & Server-exposed interface for type $\tau$ (delegates to $\mathcal{M}_{\tau}$). \\
    $r_{\tau}$ & Type-level registered resource triple $(\mathcal{C}_{\tau}, \mathcal{M}_{\tau}, \mathcal{A}_{\tau})$. \\
    \midrule
    \rowcolor{gray!15}
    \multicolumn{2}{c}{\textit{\textbf{SEPL Variables, Spaces, and Operators}}} \\
    $\mathcal{V}_{\text{evo}}$ & Universal set of evolvable variables (all managed entities plus execution artifacts). \\
    $v$ & A variable in $\mathcal{V}_{\text{evo}}$. \\
    $g_{v}$ & Learnability constraint for variable $v$ (binary). \\
    $\Theta$ & Trainable subspace, $\{v \in \mathcal{V}_{\text{evo}} \mid g_{v}=1\}$. \\
    $y$ & Execution artifacts (e.g., outputs and reasoning traces). \\
    $\mathcal{P}$ & Message space carrying auxiliary signals (traces, hypotheses, gradients, rewards) between operators. \\
    $\mathcal{P}_{\text{in}}, \mathcal{P}_{\text{out}}$ & Incoming and outgoing message types of a SEPL operator. \\
    $f$ & A SEPL operator, $f: \mathcal{V}_{\text{evo}} \times \mathcal{P}_{\text{in}} \rightarrow \mathcal{V}'_{\text{evo}} \times \mathcal{P}_{\text{out}}$. \\
    $\mathcal{Z}$ & Trace space (execution observations). \\
    $\mathcal{H}$ & Hypothesis space (causal failure attributions). \\
    $\mathcal{D}$ & Modification space (proposed resource changes). \\
    $\mathcal{G}$ & Objective specification (task goals and safety invariants). \\
    $\mathcal{S}$ & Evaluation space (performance metrics and safety status). \\
    $\rho,\sigma,\iota,\varepsilon,\kappa$ & Reflect, Select, Improve, Evaluate, and Commit operators (reflection instantiation). \\
    \midrule
    \rowcolor{gray!15}
    \multicolumn{2}{c}{\textit{\textbf{Optimization Loop (Alg.~\ref{alg:sepl-loop-appx})}}} \\
    $A$ & Agentic system. \\
    $T$ & Optimization budget (number of iterations). \\
    $t$ & Iteration index. \\
    $\mathcal{V}_{\text{evo}}^{(t)}$ & Evolvable state at iteration $t$. \\
    $\mathcal{P}^{(t)}$ & Message passed between operators at iteration $t$. \\
    $\mathcal{Z}^{(t)}$ & Observational trace at iteration $t$. \\
    $\mathcal{H}^{(t)}$ & Hypotheses at iteration $t$. \\
    $\mathcal{D}^{(t)}$ & Proposed modifications at iteration $t$. \\
    $\widetilde{\mathcal{V}}_{\text{evo}}^{(t+1)}$ & Candidate state after applying modifications. \\
    $\mathcal{S}^{(t+1)}$ & Evaluation result for the candidate state. \\
    \bottomrule
  \end{tabular}
\end{table}


\section{Code, Prompts, and Resources}
\label{appx_sec:code-prompts-resources}

All source code, agent prompts, and optimizer prompts for \projectname are organized as follows.

\paragraph{Agent and Optimizer Prompts.}
All system prompts and task-specific prompts used by the agents, as well as the prompts used by each optimizer instantiation (Reflection Optimizer, TextGrad, Reinforce++, and GRPO), are provided in the \texttt{autogenesis/} directory of the supplementary material.
The directory is structured by component: each subdirectory corresponds to a distinct agent role or optimizer module and contains the associated prompt templates.

\paragraph{Self-Evolving Code Agent Benchmark Data.}
The benchmark problems, test cases, and reference solutions for the Self-Evolving Code Agent Benchmark are provided in the \texttt{data/} directory of the supplementary material.
The dataset covers all collected LeetCode-derived problems across five programming languages (Python, C++, Java, JavaScript, and Go).
The evaluation scripts are located in \texttt{autogenesis/src/benchmark/}.

\section{Comparison with Other Protocols}
\label{appx_sec:comparison-with-other-protocols}

Table~\ref{appx_tab:tea_protocol_comparison} provides a structured protocol-level comparison between \projectname, Google A2A~\cite{google2025a2a}, and Anthropic MCP~\cite{anthropic2025}.
While A2A and MCP have standardized inter-agent communication and model-to-tool invocation respectively, both operate solely at the level of message passing and invocation, leaving internal resource states opaque and providing no primitives for lifecycle management, version lineage, or controlled state mutation. These are precisely the three properties (\textbf{Decoupling}, \textbf{Safety \& Auditability}, and \textbf{Formalism}) that \projectname is designed to provide.
The comparison is organized into five dimensions (grey rows), with blue-highlighted entries marking capabilities that are prerequisite for closed-loop self-evolution but absent from communication- or invocation-centric protocols.

\begin{table}[H]
  \centering
  \caption{Protocol-level comparison: \textbf{Autogenesis Protocol (AGP)} vs. Google A2A vs. Anthropic MCP across key dimensions for agentic systems and self-evolution.
  Symbols: \textcolor{green!60!black}{$\checkmark$} = Supported,
  \textcolor{orange}{$\triangle$} = Partial,
  \textcolor{red}{$\times$} = Not supported.
  Highlighted rows (blue background) emphasize evolution-enabling capabilities.}
  \label{appx_tab:tea_protocol_comparison}
  \footnotesize
  \renewcommand{\arraystretch}{1.15}
  \setlength{\tabcolsep}{3pt}
  \setlength{\extrarowheight}{0pt}
  \begin{tabular}{p{3cm}p{3.5cm}p{4.5cm}p{2cm}}
    \toprule
    \textbf{Dimension} & \textbf{AGP} & \textbf{A2A} & \textbf{MCP} \\
    \midrule
    \rowcolor{gray!15}
    \multicolumn{4}{c}{\textit{\textbf{Basic Information}}} \\
    \textbf{Proposer} & Our work & Google & Anthropic \\
    \textbf{Protocol Focus} & Self-evolution Agentic System & Multi-agent System Collaboration & Tool \\
    \textbf{Entity Scope} & Prompt/Agent/Tool/Env/Memory & Agent/Tool & Tool \\
    \midrule
    \rowcolor{gray!15}
    \multicolumn{4}{c}{\textit{\textbf{Agent and System Capabilities}}} \\
    \textbf{Agent First-Class} & \textcolor{green!60!black}{$\checkmark$} & \textcolor{green!60!black}{$\checkmark$} & \textcolor{red}{$\times$} \\
    \textbf{Multi-Agent} & \textcolor{green!60!black}{$\checkmark$} & \textcolor{green!60!black}{$\checkmark$} & \textcolor{red}{$\times$} \\
    \textbf{Tracer} & \textcolor{green!60!black}{$\checkmark$} & \textcolor{orange}{$\triangle$} & \textcolor{red}{$\times$} \\
    \textbf{Memory as Resource} & \textcolor{green!60!black}{$\checkmark$} & \textcolor{red}{$\times$} & \textcolor{red}{$\times$} \\
    \midrule
    \rowcolor{gray!15}
    \multicolumn{4}{c}{\textit{\textbf{Evolvable Resource Management}}} \\
    \rowcolor{blue!10}
    \textbf{Lifecycle Ops} & \textcolor{green!60!black}{$\checkmark$} & \textcolor{orange}{$\triangle$} & \textcolor{red}{$\times$} \\
    \rowcolor{blue!10}
    \textbf{Versioning and Rollback} & \textcolor{green!60!black}{$\checkmark$} & \textcolor{red}{$\times$} & \textcolor{red}{$\times$} \\
    \textbf{Registry and Retrieval} & \textcolor{green!60!black}{$\checkmark$} & \textcolor{orange}{$\triangle$} & \textcolor{orange}{$\triangle$} \\
    \rowcolor{blue!10}
    \textbf{Contract Generation} & \textcolor{green!60!black}{$\checkmark$} & \textcolor{orange}{$\triangle$} & \textcolor{red}{$\times$} \\
    \midrule
    \rowcolor{gray!15}
    \multicolumn{4}{c}{\textit{\textbf{Self-Evolution Mechanism}}} \\
    \rowcolor{blue!10}
    \textbf{Closed-Loop Evolution} & \textcolor{green!60!black}{$\checkmark$} & \textcolor{red}{$\times$} & \textcolor{red}{$\times$} \\
    \rowcolor{blue!10}
    \textbf{Operatorized Updates} & \textcolor{green!60!black}{$\checkmark$} & \textcolor{red}{$\times$} & \textcolor{red}{$\times$} \\
    \rowcolor{blue!10}
    \textbf{Auditability} & \textcolor{green!60!black}{$\checkmark$} & \textcolor{orange}{$\triangle$} & \textcolor{orange}{$\triangle$} \\
    \midrule
    \rowcolor{gray!15}
    \multicolumn{4}{c}{\textit{\textbf{General and Ecosystem}}} \\
    \textbf{Model-Agnostic} & \textcolor{green!60!black}{$\checkmark$} & \textcolor{green!60!black}{$\checkmark$} & \textcolor{green!60!black}{$\checkmark$} \\
    \textbf{Scalability} & \textcolor{green!60!black}{$O(\log n)$} & $O(n^2)$ & $O(n)$ \\
    \textbf{Open Ecosystem} & \textcolor{green!60!black}{$\checkmark$} & \textcolor{orange}{$\triangle$} & \textcolor{orange}{$\triangle$} \\
    \bottomrule
  \end{tabular}
\end{table}

\subsection{Basic Information}

\textbf{Proposer.} Google's A2A~\cite{google2025a2a} is introduced as a protocol for multi-agent communication, enabling agents to collaborate through standardized interaction primitives. Anthropic's MCP~\cite{anthropic2025} standardizes model to tool invocation interfaces. In contrast, \projectname is proposed in this work as a self-evolution protocol for composable, auditable, and safely updateable agentic systems.

\textbf{Protocol Focus.} \projectname focuses on closed-loop improvement of agentic systems by organizing resource updates through typed protocol operators and versioned state transitions. A2A primarily addresses inter-agent communication and task delegation. MCP primarily addresses standardized model to tool invocation.

\textbf{Entity Scope.} \projectname governs heterogeneous entities, including prompts, agents, tools, environments, and memory, as protocol-registered resources with explicit state and version lineage. This design supports component-level evolution, including prompt refinement and tool code updates. A2A treats agents and tools as interaction endpoints without unified lifecycle management. MCP exposes tools as callable interfaces but does not model them as evolvable components with version lineage.

\subsection{Agent and System Capabilities}

\textbf{Agent First-Class.} \projectname models agents as managed protocol components with explicit schemas, metadata, and lifecycle hooks. This enables registration, discovery, orchestration, and controlled updates. A2A is agent-centric but treats agents primarily as service endpoints without unified lifecycle management or version lineage. MCP does not define agents as protocol components and instead focuses on model to tool connectivity.

\textbf{Multi-Agent.} \projectname supports multi-agent configurations as part of its system substrate, enabling coordinated execution with traceability and evolution-ready state. A2A directly supports agent-to-agent collaboration. MCP does not treat multi-agent orchestration as a protocol-level concern.

\textbf{Execution Tracing.} \projectname provides protocol-level trace capture over inputs, outputs, intermediate decisions, and tool calls. These traces provide the learning signals required for auditable evolution. A2A and MCP leave tracing to application-level instrumentation, which can lead to inconsistent observability across deployments.

\textbf{Memory as Resource.} \projectname models memory as a first-class protocol resource with explicit read and write interfaces, state, and version lineage. This enables persistent cross-task improvement and reproducible evolution. A2A and MCP do not prescribe a memory management protocol and instead delegate persistence to external systems.

\subsection{Evolvable Resource Management}

\textbf{Lifecycle Ops.} \projectname provides standardized lifecycle operators for initialization, registration, construction, and decommissioning. These operators ensure that updates are applied to well-defined and protocol-governed targets. A2A offers partial lifecycle support for agents. MCP does not define lifecycle management across heterogeneous component types.

\textbf{Versioning and Rollback.} Version lineage and rollback form the safety foundation of closed-loop evolution. Each update produces an immutable snapshot that supports comparison, auditing, and restoration after regressions. \projectname integrates versioning as a first-class protocol capability. A2A and MCP do not natively support version lineage over protocol-managed components, which limits systematic evolution.

\textbf{Registry and Retrieval.} \projectname maintains a unified registry of protocol-registered resources and supports semantic retrieval to reduce duplication and improve composability across tasks. A2A and MCP provide partial discovery mechanisms, but they do not define a unified management plane over heterogeneous component types.

\textbf{Contract Generation.} \projectname supports automated generation of consolidated capability specifications that enumerate tool actions, arguments, preconditions, and usage constraints. This provides a principled form of context engineering that reduces prompt bloat and improves orchestration reliability. A2A and MCP rely on static descriptions or application-layer documentation without protocol-level contract aggregation.

\subsection{Self-Evolution Mechanism}

\textbf{Closed-Loop Evolution.} \projectname is built around an iterative improvement loop consisting of execution, reflection, proposal generation, evaluation, and commitment. This loop enables sustained and evidence-grounded refinement rather than one-off adaptation. A2A and MCP do not provide native self-evolution primitives.

\textbf{Operatorized Updates.} \projectname expresses state mutations as typed and composable SEPL operators with well-defined input and output contracts. This enables controlled and repeatable evolution. A2A and MCP do not define a composable operator interface for resource modification, leaving updates to application-specific logic.

\textbf{Auditability.} \projectname enforces auditability at the protocol level by recording each state transition, the execution evidence that motivated it, and the evaluation outcome that justified it. This audit trail is supported by version lineage and rollback. A2A and MCP provide only partial audit trails through external instrumentation and do not offer protocol-level guarantees.

\subsection{General and Ecosystem}

\textbf{Model-Agnostic.} This dimension assesses whether a protocol can operate across different LLM backends and providers. \projectname is model-agnostic by design through a unified model interface layer. A2A and MCP are also broadly model-agnostic because they define interaction standards rather than binding the protocol to a specific model.

\textbf{Scalability.} Scalability characterizes how coordination and discovery behave as the number of components increases. \projectname supports scalable management by treating heterogeneous components as registry-governed resources with retrieval mechanisms, enabling efficient lookup and controlled orchestration. A2A may incur increasing coordination overhead as interactions become denser in large multi-agent settings. MCP standardizes tool interfaces but may still require application-level orchestration for large tool or resource sets.

\textbf{Open Ecosystem.} Open ecosystem support refers to whether a protocol can enable reusable and interoperable components. \projectname provides a protocol stack for managing, evolving, and auditing agentic components, which supports component sharing and safe integration. A2A and MCP provide partial ecosystem support through interoperability and tool interface standardization, but they typically require additional layers for evolution-ready resource management.

\section{Details of the Self-Evolving Code Agent Benchmark}
\label{app:coding-benchmark}

\subsection{Benchmark Design Rationale}

Our benchmark is designed to evaluate self-evolving code agents under execution-grounded and human-referenced conditions. Unlike conventional code generation benchmarks that primarily assess final correctness, self-evolving agents can improve within a single inference episode by producing an initial solution, observing execution feedback, reflecting on failure modes, and revising the solution accordingly. This adaptive process requires a benchmark that measures not only whether the final submission is accepted, but also how performance evolves throughout refinement. Accordingly, our benchmark is motivated by three objectives: (i) evaluating inference-time self-evolution on executable code, (ii) calibrating agent performance against human submission distributions, and (iii) assessing cross-language robustness under long-tail language usage.

The first objective is to make self-evolution directly measurable during inference. In algorithmic coding tasks, execution feedback provides concrete and fine-grained signals, including compilation status, runtime errors, wrong answers, time-limit violations, memory-limit violations, and execution statistics for accepted submissions. These signals allow the agent to identify whether a failure stems from syntax errors, interface mismatches, corner-case logic, algorithmic inefficiency, or excessive resource usage. A benchmark for self-evolving agents should therefore expose such feedback at each refinement round and record the resulting improvement trajectory. This design distinguishes agents that solve problems through stable and efficient refinement from those that achieve correctness only through costly or unstable trial-and-error behavior.

The second objective is to evaluate coding performance relative to human submissions. Absolute pass rates are informative but insufficient for assessing practical coding competence, since they do not indicate whether an accepted solution is efficient compared with human-written solutions. We therefore build on the LeetCode online judge, which reports runtime and memory usage for accepted submissions, together with percentile-based \emph{runtime beats} and \emph{memory beats} statistics computed from human submission distributions. These human-referenced metrics provide an interpretable basis for assessing whether self-evolution improves not only correctness, but also competitiveness relative to human programmers.

The third objective is to evaluate robustness across programming languages, including long-tail languages. Many coding benchmarks are dominated by Python or other high-resource languages, which can obscure language-specific failures related to syntax, libraries, typing discipline, compilation, and runtime behavior. LeetCode provides standardized starter code across a broad set of languages, enabling the same problem to be evaluated under comparable interfaces in Python3, C++, Java, Go, and Kotlin. This design supports systematic analysis of whether self-evolution generalizes across languages and whether feedback-driven refinement remains effective across both high-resource and lower-resource programming ecosystems.

Overall, the benchmark provides a controlled setting for evaluating self-evolving code agents as dynamic problem solvers. By combining execution-based judging, iterative feedback, human-referenced efficiency statistics, and multi-language evaluation, it jointly measures functional correctness, resource efficiency, refinement dynamics, and human-relative competitiveness under a unified protocol.

\subsection{Benchmark Construction}

\textbf{Data Collection.}
We collect the full set of 3{,}822 programming problems available on LeetCode at the time of crawling. For each problem, we extract the natural-language statement, official input and output examples, constraints, platform-provided difficulty label, topical tags, and language-specific starter code templates. The topical tags characterize the algorithmic concepts required by each problem, including arrays, trees, graphs, dynamic programming, greedy methods, binary search, and mathematics. These annotations support stratified analysis across difficulty levels, algorithmic categories, and programming languages. Figure~\ref{fig:leetcode-tags-distribution} summarizes the tag and difficulty distributions of the selected problems.

\begin{figure}[H]
    \centering
    \begin{minipage}[t]{0.67\textwidth}
        \centering
        \includegraphics[width=\linewidth]{figures/evocodebench.pdf}
        \caption{Self-evolving code agent benchmark evaluation pipeline.}
        \label{fig:code-evaluation-pipeline}
    \end{minipage}\hfill
    \begin{minipage}[t]{0.31\textwidth}
        \centering
        \includegraphics[width=\linewidth]{figures/leetcode_tags_distribution.pdf}
        \caption{Problem distribution.}
        \label{fig:leetcode-tags-distribution}
    \end{minipage}
\end{figure}

The collected data are normalized into a unified problem representation. Each instance contains a fixed task specification, official examples, a language-specific starter template, and metadata for difficulty and topic categories. We preserve the original interface required by the online judge so that generated code can be submitted without modifying function signatures or class definitions. This design ensures that performance differences arise from agent behavior rather than inconsistencies in task formatting or evaluation interfaces. We conduct quality checks by filtering malformed records, removing duplicates, and verifying that starter templates are available for all target languages. From the full pool, we select 100 recently released problems to mitigate training-data contamination. The selected problems span diverse topical categories and difficulty levels, and are instantiated across Python3, C++, Java, Go, and Kotlin to enable controlled cross-language evaluation.

\textbf{Problem Characteristics.}
LeetCode-style algorithmic problems provide a controlled and challenging setting for evaluating code-agent competence. Each task specifies explicit constraints and precise input and output behavior, requiring instruction following, edge-case coverage, and faithful implementation under a fixed interface. The breadth of tags and difficulty levels evaluates algorithm selection, data-structure proficiency, and complexity-aware reasoning. Because evaluation is execution-based, brittle solutions can be exposed through concrete failures such as off-by-one errors, corner-case bugs, interface mismatches, and language-specific pitfalls. Standardized starter templates across languages further enable systematic cross-language comparison, including robustness analysis in long-tail languages. Since a solution can be revised within the same inference episode, this setting directly measures agentic capabilities such as self-repair, feedback-grounded hypothesis testing, and efficiency-aware optimization under runtime and memory constraints.

\textbf{Problem Evaluation.}
For each problem, the agent receives a fixed input representation and submits generated code to the official execution-based judge, which evaluates functional correctness on hidden test cases and reports resource usage statistics. This protocol ensures all agents are assessed under identical task inputs, execution conditions, and scoring criteria. When evaluating agents with self-evolution capability, the agent is additionally allowed to iteratively refine its solution within the same problem-solving episode under a fixed budget of 3 rounds, using execution feedback from the online judge to reflect on failure causes, identify actionable error patterns, and propose targeted code revisions, while keeping the task specification, prompt schema, and evaluation interface unchanged.

\subsection{Evaluation Metrics}

\begin{table}[H]
  \centering
  \footnotesize
  \renewcommand{\arraystretch}{0.85}
  \setlength{\tabcolsep}{4pt}
  \caption{Evaluation metrics for the algorithmic coding benchmark.}
  \label{tab:coding-metrics}
  \begin{tabular}{p{1.2cm}p{9cm}}
  \toprule
  \textbf{Metric} & \textbf{Description} \\
  \midrule
  \rowcolor{gray!15} \multicolumn{2}{c}{\textbf{Capability metrics}} \\
  \textbf{PR}  & Number of problems passing all hidden test cases within time and memory limits. \\
  \textbf{TLE} & Number of problems exceeding the allowed execution time limit. \\
  \textbf{MLE} & Number of problems exceeding the allowed memory usage. \\
  \textbf{CE}  & Number of problems where generated code failed to compile. \\
  \textbf{RE}  & Number of problems encountering a runtime error during execution. \\
  \textbf{WA}  & Number of problems producing incorrect output. \\
  \textbf{TO}  & Number of problems where the model failed to respond within the timeout. \\
  \textbf{RpE} & Number of problems where the model returned an invalid or unparseable response. \\
  \midrule
  \rowcolor{gray!15} \multicolumn{2}{c}{\textbf{Efficiency metrics}} \\
  \textbf{AR}  & Mean runtime in milliseconds over accepted solutions. \\
  \textbf{AM}  & Mean memory usage in megabytes over accepted solutions. \\
  \textbf{APC} & Mean number of test cases passed before failure. \\
  \midrule
  \rowcolor{gray!15} \multicolumn{2}{c}{\textbf{Human-referenced metrics}} \\
  \textbf{ARB} & Percentage of accepted solutions whose runtime outperforms human submissions. \\
  \textbf{AMB} & Percentage of accepted solutions whose memory usage outperforms human submissions. \\
  \bottomrule
  \end{tabular}
\end{table}

We report three groups of metrics that capture complementary aspects of code-agent performance. Capability metrics evaluate functional correctness and diagnose execution-blocking failure modes. PR measures the number of fully accepted problems, while TLE, MLE, CE, RE, WA, TO, and RpE identify whether failures arise from algorithmic inefficiency, excessive memory use, compilation errors, runtime exceptions, incorrect logic, response timeout, or invalid output formatting. These metrics are particularly important for self-evolving agents because different failure modes correspond to different refinement opportunities.

Efficiency metrics characterize the computational quality of accepted solutions and the progress made by partially correct submissions. AR and AM summarize runtime and memory usage over accepted solutions, which allows us to assess whether self-evolution improves efficiency rather than merely increasing pass rate. APC measures how many test cases are passed before failure and provides a fine-grained signal for unsuccessful submissions. This metric is useful when an agent progresses from early failure to passing most hidden tests, even if the final solution is not accepted.

Human-referenced metrics situate accepted agent solutions within the empirical distribution of human submissions. ARB measures the fraction of accepted human submissions whose runtime is slower than the agent solution, while AMB measures the analogous fraction for memory usage. These metrics provide an interpretable basis for comparing evolved agent solutions with human-written solutions and help determine whether an accepted solution is merely correct or also competitive.

For self-evolving agents, all metrics can be computed at each refinement round as well as for the final submission. This enables trajectory-level evaluation of inference-time improvement, including whether correctness increases across rounds, whether runtime and memory usage improve or degrade, and whether human-relative competitiveness changes after reflection and revision. The benchmark therefore evaluates both endpoint performance and the refinement process through which an agent reaches that endpoint.


\section{Details of Self-Evolution Protocol}

\subsection{Design Motivation}

Existing LLM-based agent systems pursue self-improvement in largely ad hoc ways: prompts, tools, and memory are tightly coupled to agent logic, updated without version control, and impossible to roll back when an update degrades performance. This architecture fragility motivates the two-layer design of \projectname. We outline the core motivating principles below.

\begin{itemize}[leftmargin=*]
    \item \textbf{Decoupling substrate from logic.} In most existing frameworks, what an agent operates over (prompts, tools, memory) and how it evolves them (optimization algorithms, feedback loops) are interleaved in a single codebase. This coupling makes it difficult to swap, reuse, or safely update individual components. \projectname separates the \emph{evolvable substrate} (RSPL) from the \emph{evolution logic} (SEPL), so that any compliant optimizer can be applied to any registered resource without modifying the component itself. This modularity is essential for principled, reproducible self-evolution.

    \item \textbf{Safety and auditability through lifecycle management.} Self-evolving agents modify their own components at runtime, which introduces risks of cascading failures or undetectable regressions. Without explicit version control and rollback support, a single bad update can silently degrade system behavior. RSPL endows every resource with versioned state and a controlled mutation interface, ensuring that each evolutionary step is traceable, reversible, and subject to explicit commit or rollback decisions. SEPL enforces that updates proceed only after formal evaluation, making every change auditable by design.

    \item \textbf{Formalism over heuristics.} Prior self-evolution approaches apply modifications heuristically, for example, by prompting a model to "improve itself" or by directly patching code, with no standardized interface governing what constitutes a valid update cycle. This informality makes it impossible to guarantee safety or reason about correctness across runs. SEPL formalizes the update cycle as a closed-loop operator interface, transforming ad hoc modifications into a rigorous protocol with well-defined pre- and post-conditions. This formalism enables algorithm-agnostic instantiation: the same operator interface supports prompt optimization, reinforcement learning, and gradient-free search.

    \item \textbf{Uniform abstraction across heterogeneous components.} Agent systems compose heterogeneous entities, including LLM instructions, external tool scripts, MCP services, and in-context memory, that are typically managed through disparate, component-specific mechanisms. \projectname provides a single, unified resource entity abstraction that encompasses all five types (Prompt, Agent, Tool, Environment, Memory), enabling SEPL to apply the same evolution operators uniformly across all components without special-casing.
\end{itemize}

Together, these principles ground the two-layer architecture in a coherent design philosophy: RSPL provides a stable, typed, versioned substrate that renders agent internals observable and controllable, while SEPL provides a safe, formal operator interface that governs how those internals are updated. The remainder of this section specifies each layer in detail.

\subsection{Layer 1: Resource Substrate Protocol Layer}

The Resource Substrate Protocol Layer (RSPL) defines the evolvable substrate as a set of protocol-registered resources with explicit state, lifecycle, and version lineage. In this paper, these resources comprise (i) \emph{instructions} (\text{Prompt}), (ii) \emph{decision policies} (\text{Agent}), (iii) \emph{actuation interfaces} (\text{Tool}), which encompass native tool scripts, MCP tools~\citep{anthropic2025}, and agent skills~\citep{anthropic2025agentskills}, (iv) \emph{task/world dynamics} (\text{Environment}), and (v) \emph{persistent state} (\text{Memory}). Crucially, resources in RSPL are \emph{passive}: they encapsulate no optimization logic and cannot self-modify; all observations and state transitions occur only through controlled, interface-mediated operations invoked by higher layers.

\subsubsection{Core Entities}

We focus on these five entity types as a minimal yet expressive substrate for agentic systems. This choice is not intended to be exhaustive, but rather to identify a common denominator across modern agent stacks and provide a uniform target space on which SEPL can operate.

\begin{definition}[Resource Entity]
\label{appx_def:resource-entity}
A resource entity of type $\tau$ and its type-level collection can be represented as:
\begin{equation}
\begin{aligned}
e_{\tau,i} &= (n_{\tau,i},\, d_{\tau,i},\, \phi_{\tau,i},\, g_{\tau,i},\, m_{\tau,i}),\\
\mathcal{E}_{\tau} &= \{\, e_{\tau,i} \mid i \in \mathcal{I}_{\tau} \,\},
\end{aligned}
\end{equation}
where $\Tau = \{\textsc{Prompt},\textsc{Agent},\textsc{Tool},\textsc{Env},\textsc{Mem}\}$ denotes the set of RSPL entity types, $\tau \in \Tau$ indexes the entity type, $\mathcal{I}_\tau$ is the index set of resource instances of type $\tau$, and $i \in \mathcal{I}_\tau$ indexes an individual instance. Here $n_{\tau,i}$ is a unique resource name, $d_{\tau,i}$ is a short description, $\phi_{\tau,i}:\mathcal{X}_{\tau} \rightarrow \mathcal{Y}_{\tau}$ is an input-to-output mapping, $g_{\tau,i} \in \{0,1\}$ is an evolvability marker, and $m_{\tau,i}$ is an auxiliary metadata dictionary.
\end{definition}

A key motivation for making prompt, tool, and memory explicit RSPL resources is \emph{decoupling}. Many agent systems package prompts, tools, and memory as internal components of an agent, which entangles agent logic with task-specific instructions and capability bundles, increasing maintenance and limiting transfer. By externalizing them as first-class, versioned resources with standardized interfaces, the same tool-calling agent policy can be paired with different prompts and tool sets, and deployed unchanged across tasks and environments.

To support resource registration, unified management, and instantiation, RSPL stores a serializable registration record for each resource instance.
\begin{definition}[Resource Registration Record]
  \label{appx_def:resource-registration-record}
  A resource registration record and its type-level collection can be represented as:
  \begin{equation}
  \begin{aligned}
  c_{\tau,i} &= (e_{\tau,i},\, v_{\tau,i},\, \eta_{\tau,i},\, \theta_{\tau,i},\, \mathcal{F}_{\tau,i}),\\
  \mathcal{C}_{\tau} &= \{\, c_{\tau,i} \mid i \in \mathcal{I}_{\tau} \,\},
  \end{aligned}
  \end{equation}
  where $\tau \in \Tau$ indexes the entity type and $i \in \mathcal{I}_\tau$ indexes an individual instance. Here $e_{\tau,i}$ is the resource entity tuple defined in \cref{appx_def:resource-entity}, $v_{\tau,i} \in \mathbb{V}$ is a version string, $\eta_{\tau,i}$ is an implementation descriptor (e.g., import path, class definition, or source-code string), $\theta_{\tau,i}$ are instantiation parameters (e.g., constructor arguments), and $\mathcal{F}_{\tau,i}$ is a set of exported representations used by LLMs to interact with the resource (e.g., function-calling schema, plain text, and structured argument schema).
  \end{definition}

\begin{definition}[Protocol-registered resource]
\label{appx_def:protocol-registered-resource}
For each entity type $\tau$, let $\mathcal{R}_\tau$ denote the type-specific registry of protocol-registered resources, and let $\mathcal{R} = \bigcup_{\tau}\mathcal{R}_\tau$ denote the global registry. RSPL binds each entity type $\tau$ to a dedicated context manager $\mathcal{M}_\tau$ and a server-exposed interface $\mathcal{A}_\tau$. We represent the type-level registered resource as
\begin{equation}
r_{\tau} = (\mathcal{C}_\tau,\; \mathcal{M}_\tau,\; \mathcal{A}_\tau),
\end{equation}
where each $c_{\tau,i} \in \mathcal{C}_\tau$ is a registration record in \cref{appx_def:resource-registration-record}. The context manager $\mathcal{M}_\tau$ maintains the collection $\mathcal{C}_\tau$, the version lineage for type $\tau$, and implements lifecycle and update operations over these records; the server-exposed interface $\mathcal{A}_\tau$ encapsulates $\mathcal{M}_\tau$ and exposes a unified external interface by delegating requests to the corresponding context-manager routines.
\end{definition}

\subsubsection{Context Manager}
\label{appx_sec:context-manager}

The context manager implements the management plane for each resource type. Beyond lifecycle control and dependency constraints, it maintains (i) an active registry of materialized resources and (ii) a versioned history for restoration. Its exported API exposes operators for lifecycle (\texttt{init}, \texttt{build}), retrieval (\texttt{list}, \texttt{get}), versioning (\texttt{update}, \texttt{restore}), execution (\texttt{run}), and serialization (\texttt{save\_to\_json}, \texttt{load\_from\_json}, \texttt{save\_contract}, \texttt{load\_contract}). The manager explicitly supports \emph{contract generation}, producing a consolidated capability and constraint specification for the managed entities, which provides stable, up-to-date descriptions that improve reliability and reduce prompt bloat, enabling systematic \emph{context engineering} via controlled prompt injection. For instance, for tools (which may be native tool scripts, MCP-connected tools~\citep{anthropic2025}, or agent skills) the contract can take a \texttt{skills.md}-style form~\citep{anthropic2025agentskills} that enumerates tool actions, arguments, preconditions, and usage constraints. The exported management interface implemented by $\mathcal{M}_\tau$ and exposed by $\mathcal{A}_\tau$ are as follows:

\begin{table}[h]
\centering
\small
\caption{Operator set of Context Manager and Server Interface.}
\label{appx_tab:context-manager-interface}
\begin{tabular}{p{2.5cm}p{10.5cm}}
\toprule
\textbf{Operator} & \textbf{Description} \\
\midrule
\rowcolor{gray!15} \multicolumn{2}{c}{\textbf{Lifecycle \& Registration}} \\
$\mathtt{init}$ & Auto discover resources and register the resource configuration to the registry. \\
$\mathtt{build}$ & Build a resource instance from code and configuration. \\
$\mathtt{register}$ & Register a new resource instance with a unique name and version. \\
$\mathtt{unregister}$ & Unregister a resource instance from the active registry and version history. \\
\midrule
\rowcolor{gray!15} \multicolumn{2}{c}{\textbf{Retrieval \& Inspection}} \\
$\mathtt{get}$ & Retrieve a resource instance by name from the active registry. \\
$\mathtt{get\_info}$ & Retrieve a resource configuration by name from the active registry. \\
$\mathtt{list}$ & List all registered resource names. \\
$\mathtt{retrieve}$ & Retrieve similar resources via semantic search when supported. \\
$\mathtt{get\_state}$ & Get the current state of a resource instance when supported. \\
\midrule
\rowcolor{gray!15} \multicolumn{2}{c}{\textbf{Versioning}} \\
$\mathtt{update}$ & Update a resource implementation and generate a new version. \\
$\mathtt{copy}$ & Duplicate a resource with an optional new name and version. \\
$\mathtt{restore}$ & Restore a specific historical version by name and version string. \\
$\mathtt{get\_variables}$ & Expose resource code/configuration as evolvable variables. \\
$\mathtt{set\_variables}$ & Update resource variables and generate a new version. \\
\midrule
\rowcolor{gray!15} \multicolumn{2}{c}{\textbf{Execution}} \\
$\mathtt{run}$ & Run a resource instance with structured input. \\
\midrule
\rowcolor{gray!15} \multicolumn{2}{c}{\textbf{Serialization}} \\
$\mathtt{save\_to\_json}$ & Serialize configurations and version history to a JSON file. \\
$\mathtt{load\_from\_json}$ & Deserialize configurations and version history from a JSON file. \\
$\mathtt{save\_contract}$ & Save the contract of a resource instance to a file. \\
$\mathtt{load\_contract}$ & Load the contract of a resource instance from a file. \\
\bottomrule
\end{tabular}
\end{table}

\subsubsection{Server Interface}
\label{appx_sec:server-interface}

The server is introduced to encapsulate the context manager's internal complexity and present a stable, simplified interface for external callers. It packages heterogeneous management routines behind a uniform set of endpoints with consistent request/response semantics, while delegating the implementation details to the context manager. This separation isolates clients from internal design changes, reduces coupling, and provides a single control plane through which the protocol mediates safe, version-aware interactions with RSPL resources.

\subsubsection{Infrastructure Services}

RSPL further includes cross-cutting services that support reliable evolution, including reproducibility, safe deployment, and versioned recovery:

\textbf{Model manager.} A unified model-API layer that standardizes calls across providers (e.g., OpenAI, Anthropic, Google, and OpenRouter, etc.), while supporting routing, fallback, and cost-aware selection to keep model access consistent as components evolve.

\textbf{Version manager.} Maintains version lineage for each resource, enabling rollback, branching, and diffing. Versions are auto-incremented identifiers (e.g., semantic versions) assigned on register or update, each referencing an immutable snapshot of the configuration record and associated artifacts for auditability and reproducibility.

\textbf{Dynamic manager.} Handles serialization and deserialization of resource configurations for persistence and transfer, enabling safe hot-swapping of resource configurations at runtime without restarting the agent system.

\textbf{Trace manager.} Captures fine-grained execution traces (inputs, outputs, intermediate decisions, tool interactions, etc.) for interpretability and debugging, and as training signals for dataset synthesis and retrospective improvement.


\subsection{Layer 2: Self-Evolution Protocol Layer}
\label{sec:appendix-self-evolution-protocol-layer}

The Self-Evolution Protocol Layer (SEPL) formalizes agentic system evolution as a generalized optimization problem over a heterogeneous state space, modeling evolutionary dynamics as a state transition function governed by a strictly typed operator algebra. By mediating all state mutations through standardized RSPL interfaces, SEPL guarantees that evolution is traceable, reversible, and safe-by-construction. While this paper focuses on the reflection-driven optimizer as the primary instantiation, the same state manipulation primitives also accommodate textual-gradient methods such as TextGrad~\citep{yuksekgonul2025optimizing} and reinforcement learning approaches such as GRPO~\citep{shao2024deepseekmath} and Reinforce++~\citep{hu2025reinforce++}.

\subsubsection{Evolvable Variables}

To transition from heuristic adaptation to a systematic evolution protocol, SEPL introduces the concept of \emph{variable lifting}: projecting discrete, heterogeneous RSPL resources (e.g., tool code, system prompts, memory modules, and environment configurations) onto a unified representation of evolvable variables. This homogenizes the interaction surface for all evolutionary operators and rigorously delineates the trainable subspace via an explicit learnability mask.

\begin{definition}[Evolvable Variable Set]
We define the universal set of evolvable variables as
\begin{equation}
\mathcal{V}_{\text{evo}} = \Bigl(\bigcup_{\tau \in \Tau} \mathcal{E}_{\tau}\Bigr) \cup \{y\},
\end{equation}
where $\mathcal{E}_{\tau}$ denotes the set of resource entities of type $\tau$ governed by RSPL, and $y$ encapsulates execution artifacts (final outputs and reasoning traces) that constitute the observational basis for retrospective optimization. Each variable $v \in \mathcal{V}_{\text{evo}}$ is associated with a binary learnability constraint $g_{v} \in \{0,1\}$, strictly defining the trainable parameter subspace
\begin{equation}
\Theta = \{v \in \mathcal{V}_{\text{evo}} \mid g_{v}=1\}.
\end{equation}
\end{definition}

The evolvability marker $g_v$ allows SEPL to operate selectively: frozen components (e.g., a fixed tool API) are excluded from the trainable subspace, while designated evolvable resources (e.g., system prompts, tool implementations) are exposed for modification. This explicit masking ensures that only intended components are mutated during evolution.

\subsubsection{Operator Algebra}

\begin{definition}[SEPL Operator]
\label{appx_def:sepl-operator}
Let $\mathcal{V}_{\text{evo}}$ be the evolvable variable set and $\mathcal{P}$ a \emph{message space} carrying auxiliary signals (e.g., traces, hypotheses, gradients, or reward signals) passed between operators. A \emph{SEPL operator} is a function
\begin{equation}
f: \mathcal{V}_{\text{evo}} \times \mathcal{P}_{\text{in}} \;\rightarrow\; \mathcal{V}'_{\text{evo}} \times \mathcal{P}_{\text{out}},
\end{equation}
where $\mathcal{P}_{\text{in}}, \mathcal{P}_{\text{out}} \subseteq \mathcal{P}$ are the incoming and outgoing message types, and $\mathcal{V}'_{\text{evo}}$ is the updated evolvable state. Operators are \emph{composable}: the output $(\mathcal{V}'_{\text{evo}}, \mathcal{P}_{\text{out}})$ of one operator serves as the input to the next, enabling the construction of an evolutionary pipeline $f_n \circ \cdots \circ f_1$. All mutations to $\mathcal{V}_{\text{evo}}$ must be routed through RSPL interfaces, ensuring every state transition is versioned, auditable, and reversible regardless of the specific optimizer instantiation.
\end{definition}

The auxiliary spaces used by operators are: trace space $\mathcal{Z}$ (execution observations), hypothesis space $\mathcal{H}$ (causal failure attributions), modification space $\mathcal{D}$ (proposed resource changes), objective specification $\mathcal{G}$ (task goals and safety invariants), and evaluation space $\mathcal{S}$ (performance metrics and safety status). The five canonical operators of the reflection instantiation are $\{\rho, \sigma, \iota, \varepsilon, \kappa\}$, corresponding to \textsc{Reflect}, \textsc{Select}, \textsc{Improve}, \textsc{Evaluate}, and \textsc{Commit}, operating over these spaces in sequence. Other instantiations (TextGrad, GRPO, Reinforce++) reuse the same operator interface but replace the internal logic of individual operators, as detailed in the method-specific subsections below.

\subsubsection{Evolutionary Loop}

Given an initial evolvable state $\mathcal{V}_{\text{evo}}^{(0)}$ and an empty message $\mathcal{P}^{(0)} = \emptyset$, the evolutionary loop at each iteration $t$ applies a sequence of operators $f_1, \ldots, f_n$ in composition:
\begin{equation}
\bigl(\mathcal{V}_{\text{evo}}^{(t+1)},\, \mathcal{P}^{(t+1)}\bigr) = (f_n \circ \cdots \circ f_1)\bigl(\mathcal{V}_{\text{evo}}^{(t)},\, \mathcal{P}^{(t)}\bigr),
\end{equation}
where each $f_i$ reads the current state and incoming messages, produces an updated state and outgoing messages consumed by $f_{i+1}$. The loop repeats until convergence or budget exhaustion. By routing all state mutations through RSPL interfaces, each transition is versioned and reversible, guaranteeing that evolution is \emph{grounded} in execution data, \emph{traceable} through versioned updates, and \emph{safe-by-construction}.

The specific operator sequence instantiated by each method determines the behavior of the loop. The reflection optimizer instantiates this loop with five operators: \textsc{Reflect} maps execution traces and current state to causal failure hypotheses, \textsc{Select} identifies target evolvable entities and generates concrete modification proposals, \textsc{Improve} applies proposals via RSPL interfaces to yield a candidate state, \textsc{Evaluate} scores the candidate against the objective and safety invariants, and \textsc{Commit} conditionally accepts or rolls back the transition. TextGrad, GRPO, and Reinforce++ reuse the same loop structure but replace the internal logic of individual operators, as detailed in the method-specific subsections below.

\subsubsection{Reflection Optimizer}

\textbf{Evolvable Variables.}
In the reflection-driven instantiation, the evolvable state is given by the lifted variable set $\mathcal{V}_{\text{evo}}$ introduced above. Concretely, $\mathcal{V}_{\text{evo}}$ includes RSPL-managed resources (e.g., prompts, tools, memories, and agent components) together with execution artifacts (e.g., the produced answer and reasoning trace). A binary learnability mask specifies which variables may be modified, allowing the optimizer to target only authorized components while keeping non-learnable resources fixed.

\textbf{Operator Algebra.}
We instantiate SEPL with the canonical reflection-driven operator suite. The operator signatures and their intended roles are as follows.

\begin{itemize}[leftmargin=1em, itemsep=0em, parsep=0em, topsep=0em, partopsep=0em]
  \item \textbf{Reflect ($\rho$).} Defined as $\rho: \mathcal{Z} \times \mathcal{V}_{\text{evo}} \rightarrow \wp(\mathcal{H})$, this operator bridges the gap between raw observation and optimization direction. It approximates the ``semantic gradient'' of the system by mapping high-dimensional execution traces to specific, causal failure hypotheses within the variable space.
  \item \textbf{Select ($\sigma$).} Formulated as $\sigma: \mathcal{V}_{\text{evo}} \times \wp(\mathcal{H}) \rightarrow \wp(\mathcal{D})$, this operator acts as the targeting policy. It identifies which evolvable entities within $\mathcal{V}_{\text{evo}}$ are implicated by the diagnostic hypotheses, then generates concrete modification proposals $\mathcal{D}$ targeting those entities, subject to structural constraints.
  \item \textbf{Improve ($\iota$).} The mutation operator, $\iota: \mathcal{V}_{\text{evo}} \times \wp(\mathcal{D}) \rightarrow \mathcal{V}'_{\text{evo}}$, executes the physical state transition. It applies discrete updates $\mathcal{D}$ via standardized RSPL interfaces to yield a provisional candidate state.
  \item \textbf{Evaluate ($\varepsilon$).} Specified as $\varepsilon: \mathcal{V}'_{\text{evo}} \times \mathcal{G} \rightarrow \mathcal{S}$, this operator serves as the objective function. It maps the candidate state and goal specification to the evaluation space $\mathcal{S}$ (comprising quantitative scores and strict safety invariants).
  \item \textbf{Commit ($\kappa$).} Operating as $\kappa: \mathcal{V}'_{\text{evo}} \times \mathcal{S} \rightarrow \mathcal{V}_{\text{evo}}$, this function acts as a conditional gating mechanism. It utilizes the evaluation signals in $\mathcal{S}$ to govern state transition, rigorously enforcing safety invariants and performance monotonicity by accepting the candidate $\mathcal{V}'_{\text{evo}}$ only when specific success criteria are met.
\end{itemize}

\textbf{The Evolutionary Loop.}
These operators are composed into the reflection-driven closed-loop procedure shown in Algorithm~\ref{alg:sepl-loop-appx}. Starting from an initial lifted state $\mathcal{V}_{\text{evo}}^{(0)}$, the agent first executes to collect an observational trace $\mathcal{Z}$ (tool outputs, intermediate decisions, failures, and progress signals). The reflect operator $\rho$ maps $\mathcal{Z}$ to a set of causal hypotheses $\mathcal{H}$, which are then translated by $\sigma$ into concrete modification primitives $\mathcal{D}$ (e.g., prompt edits, tool adjustments, or memory updates) over the learnable subset of $\mathcal{V}_{\text{evo}}$. The improve operator $\iota$ applies $\mathcal{D}$ via RSPL interfaces to obtain a candidate state, which is evaluated by $\varepsilon$ to produce $\mathcal{S}$ capturing both performance metrics and safety constraints. Finally, the commit operator $\kappa$ gates the transition by accepting only candidates that satisfy the predefined criteria, recording each accepted change as a versioned resource update with auditable lineage and enabling rollback when necessary.

\begin{algorithm}[t]
  \caption{Reflection Optimizer Evolutionary Loop}
  \label{alg:sepl-loop-appx}
  % Right-aligned comments with a fixed-width comment column to avoid large blank gaps.
  \newcommand{\algcomment}[1]{\hfill\parbox[t]{0.48\linewidth}{\raggedleft\footnotesize\textcolor{gray}{$\rhd$ \textit{#1}}}}
  \begin{algorithmic}[1]
  \renewcommand{\algorithmicrequire}{\textbf{Input:}}
  \renewcommand{\algorithmicensure}{\textbf{Output:}}
  \REQUIRE Agentic System $\mathcal{A}$, Objective $\mathcal{G}$, Budget $T$
  \ENSURE Optimized state $\mathcal{V}_{\text{evo}}^{*}$

  \STATE \textbf{Initialization:}
  \STATE $\mathcal{V}_{\text{evo}}^{(0)} \leftarrow \text{VariableLifting}(\mathcal{A})$ \algcomment{Project resources to optimization manifold}
  \STATE $\mathcal{Z}^{(0)} \leftarrow \text{Execute}(\mathcal{A}, \mathcal{V}_{\text{evo}}^{(0)})$ \algcomment{Trace: tool I/O, failures, latencies, progress}
  
  \vspace{0.3em}
  \STATE \textbf{Optimization Cycle:}
  \FOR{$t = 0, 1, \ldots, T-1$}
      \STATE \textcolor{gray}{// Phase 1: Diagnosis \& Proposal}
      \STATE $\mathcal{H}^{(t)} \leftarrow \rho(\mathcal{Z}^{(t)}, \mathcal{V}_{\text{evo}}^{(t)})$ \algcomment{Reflect: attribute failures / inefficiencies}
      \STATE $\mathcal{D}^{(t)} \leftarrow \sigma(\mathcal{V}_{\text{evo}}^{(t)}, \mathcal{H}^{(t)})$ \algcomment{Select: propose edits over learnable variables}
      
      \vspace{0.2em}
      \STATE \textcolor{gray}{// Phase 2: Mutation \& Verification}
      \STATE $\widetilde{\mathcal{V}}_{\text{evo}}^{(t+1)} \leftarrow \iota(\mathcal{V}_{\text{evo}}^{(t)}, \mathcal{D}^{(t)})$ \algcomment{Improve: apply proposed updates (candidate)}
      \STATE $\mathcal{S}^{(t+1)} \leftarrow \varepsilon(\widetilde{\mathcal{V}}_{\text{evo}}^{(t+1)}, \mathcal{G})$ \algcomment{Evaluate: metrics + safety invariants}
      
      \vspace{0.2em}
      \STATE \textcolor{gray}{// Phase 3: Gating \& Transition}
      \IF{$\text{Accept}(\mathcal{S}^{(t+1)})$}
          \STATE \textcolor{gray}{// Accept: safe \& non-degrading}
          \STATE $\mathcal{V}_{\text{evo}}^{(t+1)} \leftarrow \kappa(\widetilde{\mathcal{V}}_{\text{evo}}^{(t+1)}, \mathcal{S}^{(t+1)})$ \algcomment{Commit: versioned update}
      \ELSE
          \STATE \textcolor{gray}{// Reject: rollback / keep previous state}
          \STATE $\mathcal{V}_{\text{evo}}^{(t+1)} \leftarrow \mathcal{V}_{\text{evo}}^{(t)}$
      \ENDIF
      
      \vspace{0.2em}
      \STATE \textcolor{gray}{// Phase 4: Next Iteration}
      \STATE $\mathcal{Z}^{(t+1)} \leftarrow \text{Execute}(\mathcal{A}, \mathcal{V}_{\text{evo}}^{(t+1)})$ \algcomment{Re-run under updated resources}
      \IF{$\text{Converged}(\mathcal{S}^{(t+1)})$}
          \STATE \textbf{break}
      \ENDIF
  \ENDFOR
  \STATE \textbf{return} $\mathcal{V}_{\text{evo}}^{(t)}$
  \end{algorithmic}
  \end{algorithm}

\subsubsection{TextGrad Optimizer}

\textbf{Evolvable Variables.}
In the TextGrad instantiation, the evolvable variables are restricted to a subset of \emph{prompt variables} marked as optimizable and lifted into TextGrad variables with explicit role descriptions. In our implementation, each optimizable prompt module is represented as a TextGrad variable whose value is the current prompt text and whose role description specifies the prompt's function, enabling the optimizer to condition updates on its intended semantics.

\textbf{Operator Algebra.}
TextGrad instantiates SEPL with a prompt-level operatorization in which ``gradients'' are natural-language critiques produced by an LLM evaluator and updates are implemented as constrained prompt rewrites.
Following the standard TextGrad view, we express the method with five core operators, namely \emph{Execute}, \emph{Loss}, \emph{Backward}, \emph{Improve}, and \emph{Commit}, where the ``gradient'' is a piece of text (a critique) rather than a numeric vector:
\begin{itemize}[leftmargin=1em, itemsep=0em, parsep=0em, topsep=0em, partopsep=0em]
  \item \textbf{Execute ($\chi_{\mathrm{tg}}$).} $\chi_{\mathrm{tg}}: (A, \mathcal{V}_{\text{evo}}, x, f) \rightarrow \mathcal{Z}$ runs the agent under the current prompt variables and produces an execution trace/outcome.
  \item \textbf{Loss ($\lambda_{\mathrm{tg}}$).} $\lambda_{\mathrm{tg}}: \mathcal{Z} \rightarrow \mathcal{G}_{\mathrm{tg}}$, where $\mathcal{G}_{\mathrm{tg}}$ is a space of natural-language critiques (textual gradients). In our implementation, $\lambda_{\mathrm{tg}}$ is realized by \texttt{TextLoss}, which queries an evaluator LLM and returns critique feedback.
  \item \textbf{Backward ($\beta_{\mathrm{tg}}$).} $\beta_{\mathrm{tg}}: \mathcal{V}_{\text{evo}} \times \mathcal{G}_{\mathrm{tg}} \rightarrow \mathcal{V}_{\text{evo}}$ assigns textual gradients to optimizable prompt variables by storing the critique (optionally with context) in a per-variable gradient buffer. In our current implementation, we distribute the same critique to each optimizable prompt variable for stability.
  \item \textbf{Improve ($\iota_{\mathrm{tg}}$).} $\iota_{\mathrm{tg}}: \mathcal{V}_{\text{evo}} \rightarrow \mathcal{V}'_{\text{evo}}$ rewrites prompt variables via a textual-gradient-descent step: it constructs an update instruction from each variable's role description, current value, and accumulated textual gradients, then queries an optimizer LLM and extracts the improved variable text from a constrained output format.
  \item \textbf{Commit ($\kappa_{\mathrm{tg}}$).} $\kappa_{\mathrm{tg}}: \mathcal{V}'_{\text{evo}} \rightarrow \mathcal{V}_{\text{evo}}$ synchronizes the updated prompt variables back into the running agent and clears caches, completing the state transition.
\end{itemize}

\textbf{The Evolutionary Loop.}
Algorithm~\ref{alg:textgrad-loop-appx} presents the full TextGrad optimization cycle in operator form. At each iteration, the agent is executed under the current prompt variables to obtain a trace $\mathcal{Z}$ via $\chi_{\mathrm{tg}}$, an LLM-based evaluator produces a natural-language critique $g \in \mathcal{G}_{\mathrm{tg}}$ via $\lambda_{\mathrm{tg}}$, the critique is assigned as a \emph{textual gradient} to the optimizable prompt variables via $\beta_{\mathrm{tg}}$, the prompt variables are improved via $\iota_{\mathrm{tg}}$ using textual-gradient-descent, and the candidate state is committed via $\kappa_{\mathrm{tg}}$ to synchronize the updated prompts back into the running agent (and clear caches) before the next iteration.

\begin{algorithm}[t]
  \caption{TextGrad Prompt Optimization Loop}
  \label{alg:textgrad-loop-appx}
  % Right-aligned comments with a fixed-width comment column (same style as Algorithm~\ref{alg:sepl-loop-appx}).
  \newcommand{\algcommenttg}[1]{\hfill\parbox[t]{0.48\linewidth}{\raggedleft\footnotesize\textcolor{gray}{$\rhd$ \textit{#1}}}}
  \begin{algorithmic}[1]
  \renewcommand{\algorithmicrequire}{\textbf{Input:}}
  \renewcommand{\algorithmicensure}{\textbf{Output:}}
  \REQUIRE Agentic System $\mathcal{A}$, task $x$, attachments $f$ (optional), Budget $K$, evaluator/optimizer LLMs $M_{\text{eval}}, M_{\text{opt}}$
  \ENSURE Updated state $\mathcal{V}_{\text{evo}}^{*}$ (prompt variables updated via TextGrad)

  \STATE \textcolor{gray}{// Phase 0: Setup}
  \STATE Set backward engine to $M_{\text{eval}}$ \algcommenttg{Evaluator used by TextLoss}
  \STATE $\mathcal{V}_{\text{evo}}^{(0)} \leftarrow \text{VariableLifting}(\mathcal{A})$ \algcommenttg{Lift optimizable prompts to TextGrad variables}
  \STATE Initialize textual optimizer with $M_{\text{opt}}$ \algcommenttg{TextualGradientDescent over prompt vars}

  \vspace{0.2em}
  \STATE \textcolor{gray}{// Optimization Cycle}
  \FOR{$k = 0, 1, \ldots, K-1$}
      \STATE \textcolor{gray}{// Phase 1: Execute (Forward)}
      \STATE $\mathcal{Z}^{(k)} \leftarrow \chi_{\mathrm{tg}}(\mathcal{A}, \mathcal{V}_{\text{evo}}^{(k)}, x, f)$ \algcommenttg{Run agent with current prompts}
      
      \vspace{0.1em}
      \STATE \textcolor{gray}{// Phase 2: Loss (Textual Gradient)}
      \STATE Build evaluation instruction from $\mathcal{Z}^{(k)}$ \algcommenttg{Condition on success/error}
      \STATE $g^{(k)} \leftarrow \lambda_{\mathrm{tg}}(\mathcal{Z}^{(k)})$ \algcommenttg{TextLoss produces critique string}
      
      \vspace{0.1em}
      \STATE \textcolor{gray}{// Phase 3: Backward (Assign Gradients)}
      \STATE $\mathcal{V}_{\text{evo}}^{(k)} \leftarrow \beta_{\mathrm{tg}}(\mathcal{V}_{\text{evo}}^{(k)}, g^{(k)})$ \algcommenttg{Assign critique to gradient buffers}
      
      \vspace{0.1em}
      \STATE \textcolor{gray}{// Phase 4: Improve (Textual Gradient Descent)}
      \STATE $\widetilde{\mathcal{V}}_{\text{evo}}^{(k+1)} \leftarrow \iota_{\mathrm{tg}}(\mathcal{V}_{\text{evo}}^{(k)})$ \algcommenttg{Rewrite prompts via textual GD}
      
      \vspace{0.1em}
      \STATE \textcolor{gray}{// Phase 5: Commit \& Next Iteration}
      \STATE $\mathcal{V}_{\text{evo}}^{(k+1)} \leftarrow \kappa_{\mathrm{tg}}(\widetilde{\mathcal{V}}_{\text{evo}}^{(k+1)})$ \algcommenttg{Sync back; clear caches}
      \IF{$\text{Converged}(g^{(k)})$}
          \STATE \textbf{break}
      \ENDIF
  \ENDFOR
  \STATE \textbf{return} $\mathcal{V}_{\text{evo}}^{(k)}$
  \end{algorithmic}
  \end{algorithm}

\subsubsection{Reinforce++ Optimizer}

\textbf{Evolvable Variables.}
Reinforce++ optimizes a trainable subset of RSPL resources, focusing on prompt variables and tool implementations (native scripts, MCP tools~\citep{anthropic2025}, and agent skills~\citep{anthropic2025agentskills}), and optionally refining the produced solution text. Our implementation follows a two stage structure: (i) update trainable variables that govern behavior (e.g., prompts and tools), and (ii) update the solution itself when enabled.

\textbf{Operator Algebra.}
Reinforce++ is characterized by a clipped objective with an explicit penalty to a reference solution, while using reflection to translate RL signals into concrete edits. We group the method into a small set of core operators:
\begin{itemize}[leftmargin=1em, itemsep=0em, parsep=0em, topsep=0em, partopsep=0em]
  \item \textbf{Sample ($\chi_{\mathrm{rpp}}$).} $\chi_{\mathrm{rpp}}: (A, \mathcal{V}_{\text{evo}}, x, f) \rightarrow \mathcal{Z}$ samples a rollout under the current resources and yields an execution trace containing the produced answer.
  \item \textbf{Reward ($\varepsilon_{\mathrm{rpp}}$).} $\varepsilon_{\mathrm{rpp}}: (y^{(t)}, y^{(t-1)}, y^{*}, y_{\mathrm{sft}}) \rightarrow (r^{(t)}, A^{(t)}, J^{(t)}, \pi^{(t)})$ computes the RL signal tuple from the current solution $y^{(t)}$. Here $r^{(t)}$ is a task reward comparing $y^{(t)}$ with $y^{*}$, and $\pi^{(t)}$ is a policy ratio surrogate approximated via text similarity $\eta(\cdot,\cdot)$ as $\pi^{(t)} \triangleq \eta(y^{(t-1)}, y^{(t)})$ (since token-level probability ratios are unavailable in inference-only LLM settings). We define a penalty to a reference solution $y_{\mathrm{sft}}$ as $\mathrm{pen}^{(t)} \triangleq \beta\,\bigl|\log \max(\eta(y_{\mathrm{sft}}, y^{(t)}), \epsilon_0)\bigr|$ and set $A^{(t)} \triangleq r^{(t)} - \mathrm{pen}^{(t)}$. The clipped Reinforce++ objective is
  \[
    J^{(t)} \triangleq \min\bigl(\pi^{(t)} A^{(t)},\; \bar{\pi}^{(t)} A^{(t)}\bigr),
    \quad \bar{\pi}^{(t)} \triangleq \mathrm{clip}(\pi^{(t)}, 1-\epsilon, 1+\epsilon).
  \]
  \item \textbf{Diagnose ($\delta_{\mathrm{rpp}}$).} $\delta_{\mathrm{rpp}}: (\mathcal{Z}, \mathcal{V}_{\text{train}}, r^{(t)}, A^{(t)}, J^{(t)}, \pi^{(t)}) \rightarrow \mathcal{H}$ produces an edit oriented diagnosis that is explicitly conditioned on the RL metrics and the execution trace.
  \item \textbf{Improve ($\iota_{\mathrm{rpp}}$).} $\iota_{\mathrm{rpp}}: (\mathcal{V}, \mathcal{H}) \rightarrow \mathcal{V}'_{\text{evo}}$ applies RL informed edits to either (i) the trainable resources $\mathcal{V}_{\text{train}}$ such as prompts and tools, or (ii) the solution variable itself when solution refinement is enabled, yielding a candidate state.
  \item \textbf{Commit ($\kappa_{\mathrm{rpp}}$).} $\kappa_{\mathrm{rpp}}: \mathcal{V}'_{\text{evo}} \rightarrow \mathcal{V}_{\text{evo}}$ applies accepted updates back to RSPL resources, completing the state transition.
\end{itemize}

\textbf{The Evolutionary Loop.}
Algorithm~\ref{alg:reinforcepp-loop-appx} summarizes the Reinforce++ loop in a phased form. Each iteration (i) computes Reinforce++ signals via the clipped objective and the penalty to the reference solution, (ii) improves trainable resources through RL conditioned reflection and edits, (iii) optionally improves the solution text, and (iv) applies an early stopping evaluation.

\begin{algorithm}[t]
  \caption{Reinforce++ Optimization Loop}
  \label{alg:reinforcepp-loop-appx}
  % Right-aligned comments with a fixed-width comment column.
  \newcommand{\algcommentrpp}[1]{\hfill\parbox[t]{0.48\linewidth}{\raggedleft\footnotesize\textcolor{gray}{$\rhd$ \textit{#1}}}}
  \begin{algorithmic}[1]
  \renewcommand{\algorithmicrequire}{\textbf{Input:}}
  \renewcommand{\algorithmicensure}{\textbf{Output:}}
  \REQUIRE Agentic System $\mathcal{A}$, task $x$, attachments $f$ (optional), ground truth $y^{*}$, reference solution $y_{\mathrm{sft}}$, Budget $T$
  \ENSURE Final solution $y^{(t)}$ and updated trainable resources $\mathcal{V}_{\text{train}}$

  \STATE \textcolor{gray}{// Initialization}
  \STATE $\mathcal{V}_{\text{evo}}^{(0)} \leftarrow \text{VariableLifting}(\mathcal{A})$ \algcommentrpp{Lift trainable resources}
  \STATE $\mathcal{Z}^{(0)} \leftarrow \chi_{\mathrm{rpp}}(\mathcal{A}, \mathcal{V}_{\text{evo}}^{(0)}, x, f)$ \algcommentrpp{Sample once}
  \STATE Extract solution $y^{(0)}$ from $\mathcal{Z}^{(0)}$
  \STATE $y^{(-1)} \leftarrow y^{(0)}$ \algcommentrpp{Initialize previous solution}
  
  \FOR{$t = 0, 1, \ldots, T-1$}
      \STATE \textcolor{gray}{// Phase 1: Reinforce++ reward and objective}
      \STATE $(r^{(t)}, A^{(t)}, J^{(t)}, \pi^{(t)}) \leftarrow \varepsilon_{\mathrm{rpp}}(y^{(t)}, y^{(t-1)}, y^{*}, y_{\mathrm{sft}})$ \algcommentrpp{Reward, penalty, clipped objective}

      \vspace{0.1em}
      \STATE \textcolor{gray}{// Phase 2: Improve trainable resources (prompt and tool)}
      \STATE $\mathcal{V}_{\text{train}}^{(t)} \leftarrow \text{GetTrainables}(\mathcal{V}_{\text{evo}}^{(t)})$
      \STATE $\mathcal{H}_{\text{train}}^{(t)} \leftarrow \delta_{\mathrm{rpp}}(\mathcal{Z}^{(t)}, \mathcal{V}_{\text{train}}^{(t)}, r^{(t)}, A^{(t)}, J^{(t)}, \pi^{(t)})$ \algcommentrpp{Diagnose conditioned on RL signals}
      \STATE $\widetilde{\mathcal{V}}_{\text{train}}^{(t+1)} \leftarrow \iota_{\mathrm{rpp}}(\mathcal{V}_{\text{train}}^{(t)}, \mathcal{H}_{\text{train}}^{(t)})$ \algcommentrpp{Apply edits to trainables (candidate)}
      \STATE $\mathcal{V}_{\text{train}}^{(t+1)} \leftarrow \kappa_{\mathrm{rpp}}(\widetilde{\mathcal{V}}_{\text{train}}^{(t+1)})$ \algcommentrpp{Commit updates}

      \vspace{0.1em}
      \STATE \textcolor{gray}{// Phase 3: Re run under updated resources}
      \STATE $\mathcal{Z}^{(t+1)} \leftarrow \chi_{\mathrm{rpp}}(\mathcal{A}, \mathcal{V}_{\text{evo}}^{(t)} \cup \mathcal{V}_{\text{train}}^{(t+1)}, x, f)$
      \STATE Extract solution $y^{(t+1)}$ from $\mathcal{Z}^{(t+1)}$

      \vspace{0.1em}
      \STATE \textcolor{gray}{// Phase 4: Optional solution refinement}
      \STATE $\mathcal{H}_{\text{sol}}^{(t)} \leftarrow \delta_{\mathrm{rpp}}(\mathcal{Z}^{(t+1)}, \{y^{(t+1)}\}, r^{(t)}, A^{(t)}, J^{(t)}, \pi^{(t)})$ \algcommentrpp{Diagnose solution quality}
      \STATE $\widetilde{y}^{(t+1)} \leftarrow \iota_{\mathrm{rpp}}(y^{(t+1)}, \mathcal{H}_{\text{sol}}^{(t)})$ \algcommentrpp{Edit solution text (candidate)}
      \STATE $y^{(t+1)} \leftarrow \kappa_{\mathrm{rpp}}(\widetilde{y}^{(t+1)})$ \algcommentrpp{Commit solution update}
      
      \vspace{0.1em}
      \STATE \textcolor{gray}{// Phase 5: Early stopping}
      \IF{$\text{Satisfied}(\mathcal{Z}^{(t+1)})$}
          \STATE \textbf{break}
      \ENDIF
      \STATE $y^{(t)} \leftarrow y^{(t+1)}$ \algcommentrpp{Advance current solution}
  \ENDFOR
  \STATE \textbf{return} $y^{(t)}$
  \end{algorithmic}
  \end{algorithm}

\subsubsection{GRPO Optimizer}

\textbf{Evolvable Variables.}
GRPO optimizes a trainable subset of RSPL resources, focusing on prompt variables and tool implementations (native scripts, MCP tools~\citep{anthropic2025}, and agent skills~\citep{anthropic2025agentskills}), and optionally refining the produced solution text. Similar to Reinforce++, our implementation follows a two stage structure: (i) update trainable variables that govern behavior (e.g., prompts and tools), and (ii) update the solution itself when enabled.

\textbf{Operator Algebra.}
GRPO is characterized by sampling multiple candidate solutions per step and using group normalized advantages with a clipped objective. We formalize the method with the following core operators:
\begin{itemize}[leftmargin=1em, itemsep=0em, parsep=0em, topsep=0em, partopsep=0em]
  \item \textbf{Sample ($\chi_{\mathrm{grpo}}$).} $\chi_{\mathrm{grpo}}: (A, \mathcal{V}_{\text{evo}}, x, f, K) \rightarrow \{\mathcal{Z}_i\}_{i=1}^K$ samples $K$ independent rollouts under the current resources, yielding $K$ execution traces each containing a candidate solution $y_i$.
  \item \textbf{Reward ($\varepsilon_{\mathrm{grpo}}$).} $\varepsilon_{\mathrm{grpo}}: (\{y_i\}_{i=1}^K, y^{*}, y^{(t-1)}) \rightarrow (\{r_i\}_{i=1}^K, \{A_i\}_{i=1}^K, \{J_i\}_{i=1}^K, \{\pi_i\}_{i=1}^K)$ computes RL signals for all $K$ candidates. For each candidate $y_i$, we compute a task reward $r_i$ comparing $y_i$ with $y^{*}$, a policy ratio surrogate $\pi_i \triangleq \eta(y^{(t-1)}, y_i)$ approximated via text similarity $\eta(\cdot,\cdot)$ (since token-level probability ratios are unavailable in inference-only LLM settings), and a group normalized advantage $A_i$ by normalizing rewards across the candidate set: $A_i = (r_i - \bar{r}) / \sigma_r$ where $\bar{r}$ and $\sigma_r$ are the mean and standard deviation of $\{r_i\}_{i=1}^K$. The GRPO clipped objective for each candidate is
  \[
    J_i \triangleq \min\bigl(\pi_i A_i,\; \bar{\pi}_i A_i\bigr),
    \quad \bar{\pi}_i \triangleq \begin{cases}
      \min(\pi_i, 1+\epsilon) & \text{if } A_i \geq 0 \\
      \max(\pi_i, 1-\epsilon) & \text{if } A_i < 0
    \end{cases}.
  \]
  \item \textbf{Diagnose ($\delta_{\mathrm{grpo}}$).} $\delta_{\mathrm{grpo}}: (\{\mathcal{Z}_i\}_{i=1}^K, \mathcal{V}_{\text{train}}, \{r_i, A_i, J_i, \pi_i\}_{i=1}^K) \rightarrow \mathcal{H}$ produces an edit oriented diagnosis that is explicitly conditioned on the multiple candidate solutions and their RL metrics, enabling the optimizer to identify patterns across candidates.
  \item \textbf{Improve ($\iota_{\mathrm{grpo}}$).} $\iota_{\mathrm{grpo}}: (\mathcal{V}, \mathcal{H}) \rightarrow \mathcal{V}'_{\text{evo}}$ applies RL informed edits to either (i) the trainable resources $\mathcal{V}_{\text{train}}$ such as prompts and tools, or (ii) the solution variable itself when solution refinement is enabled, yielding a candidate state.
  \item \textbf{Commit ($\kappa_{\mathrm{grpo}}$).} $\kappa_{\mathrm{grpo}}: \mathcal{V}'_{\text{evo}} \rightarrow \mathcal{V}_{\text{evo}}$ applies accepted updates back to RSPL resources, completing the state transition.
\end{itemize}

\textbf{The Evolutionary Loop.}
Algorithm~\ref{alg:grpo-loop-appx} summarizes the GRPO loop in a phased form. Each iteration (i) samples $K$ candidate solutions, (ii) computes GRPO signals via group normalized advantages and clipped objectives, (iii) improves trainable resources through multi candidate conditioned reflection and edits, (iv) optionally improves the solution text, and (v) applies an early stopping evaluation.

\begin{algorithm}[t]
  \caption{GRPO Optimization Loop}
  \label{alg:grpo-loop-appx}
  % Right-aligned comments with a fixed-width comment column.
  \newcommand{\algcommentgrpo}[1]{\hfill\parbox[t]{0.48\linewidth}{\raggedleft\footnotesize\textcolor{gray}{$\rhd$ \textit{#1}}}}
  \begin{algorithmic}[1]
  \renewcommand{\algorithmicrequire}{\textbf{Input:}}
  \renewcommand{\algorithmicensure}{\textbf{Output:}}
  \REQUIRE Agentic System $\mathcal{A}$, task $x$, attachments $f$ (optional), ground truth $y^{*}$, Budget $T$, number of candidates $K$
  \ENSURE Final solution $y^{(t)}$ and updated trainable resources $\mathcal{V}_{\text{train}}$

  \STATE \textcolor{gray}{// Initialization}
  \STATE $\mathcal{V}_{\text{evo}}^{(0)} \leftarrow \text{VariableLifting}(\mathcal{A})$ \algcommentgrpo{Lift trainable resources}
  \STATE $\mathcal{Z}^{(0)} \leftarrow \chi_{\mathrm{grpo}}(\mathcal{A}, \mathcal{V}_{\text{evo}}^{(0)}, x, f, 1)$ \algcommentgrpo{Sample initial solution}
  \STATE Extract solution $y^{(0)}$ from $\mathcal{Z}^{(0)}$
  \STATE $y^{(-1)} \leftarrow y^{(0)}$ \algcommentgrpo{Initialize previous solution}
  
  \FOR{$t = 0, 1, \ldots, T-1$}
      \STATE \textcolor{gray}{// Phase 1: Sample multiple candidates}
      \STATE $\{\mathcal{Z}_i^{(t)}\}_{i=1}^K \leftarrow \chi_{\mathrm{grpo}}(\mathcal{A}, \mathcal{V}_{\text{evo}}^{(t)}, x, f, K)$ \algcommentgrpo{Sample $K$ rollouts}
      \STATE Extract candidate solutions $\{y_i^{(t)}\}_{i=1}^K$ from $\{\mathcal{Z}_i^{(t)}\}_{i=1}^K$
      
      \vspace{0.1em}
      \STATE \textcolor{gray}{// Phase 2: GRPO reward and objective}
      \STATE $(\{r_i^{(t)}\}_{i=1}^K, \{A_i^{(t)}\}_{i=1}^K, \{J_i^{(t)}\}_{i=1}^K, \{\pi_i^{(t)}\}_{i=1}^K) \leftarrow \varepsilon_{\mathrm{grpo}}(\{y_i^{(t)}\}_{i=1}^K, y^{*}, y^{(t-1)})$ \algcommentgrpo{Group normalized advantages, clipped objectives}

      \vspace{0.1em}
      \STATE \textcolor{gray}{// Phase 3: Improve trainable resources (prompt and tool)}
      \STATE $\mathcal{V}_{\text{train}}^{(t)} \leftarrow \text{GetTrainables}(\mathcal{V}_{\text{evo}}^{(t)})$
      \STATE $\mathcal{H}_{\text{train}}^{(t)} \leftarrow \delta_{\mathrm{grpo}}(\{\mathcal{Z}_i^{(t)}\}_{i=1}^K, \mathcal{V}_{\text{train}}^{(t)}, \{r_i^{(t)}, A_i^{(t)}, J_i^{(t)}, \pi_i^{(t)}\}_{i=1}^K)$ \algcommentgrpo{Diagnose conditioned on multi candidate RL signals}
      \STATE $\widetilde{\mathcal{V}}_{\text{train}}^{(t+1)} \leftarrow \iota_{\mathrm{grpo}}(\mathcal{V}_{\text{train}}^{(t)}, \mathcal{H}_{\text{train}}^{(t)})$ \algcommentgrpo{Apply edits to trainables (candidate)}
      \STATE $\mathcal{V}_{\text{train}}^{(t+1)} \leftarrow \kappa_{\mathrm{grpo}}(\widetilde{\mathcal{V}}_{\text{train}}^{(t+1)})$ \algcommentgrpo{Commit updates}

      \vspace{0.1em}
      \STATE \textcolor{gray}{// Phase 4: Re run under updated resources}
      \STATE $\mathcal{Z}^{(t+1)} \leftarrow \chi_{\mathrm{grpo}}(\mathcal{A}, \mathcal{V}_{\text{evo}}^{(t)} \cup \mathcal{V}_{\text{train}}^{(t+1)}, x, f, 1)$
      \STATE Extract solution $y^{(t+1)}$ from $\mathcal{Z}^{(t+1)}$

      \vspace{0.1em}
      \STATE \textcolor{gray}{// Phase 5: Optional solution refinement}
      \STATE $\mathcal{H}_{\text{sol}}^{(t)} \leftarrow \delta_{\mathrm{grpo}}(\{\mathcal{Z}_i^{(t)}\}_{i=1}^K, \{y^{(t+1)}\}, \{r_i^{(t)}, A_i^{(t)}, J_i^{(t)}, \pi_i^{(t)}\}_{i=1}^K)$ \algcommentgrpo{Diagnose solution quality using multi candidate context}
      \STATE $\widetilde{y}^{(t+1)} \leftarrow \iota_{\mathrm{grpo}}(y^{(t+1)}, \mathcal{H}_{\text{sol}}^{(t)})$ \algcommentgrpo{Edit solution text (candidate)}
      \STATE $y^{(t+1)} \leftarrow \kappa_{\mathrm{grpo}}(\widetilde{y}^{(t+1)})$ \algcommentgrpo{Commit solution update}
      
      \vspace{0.1em}
      \STATE \textcolor{gray}{// Phase 6: Early stopping}
      \IF{$\text{Satisfied}(\mathcal{Z}^{(t+1)})$}
          \STATE \textbf{break}
      \ENDIF
      \STATE $y^{(t)} \leftarrow y^{(t+1)}$ \algcommentgrpo{Advance current solution}
  \ENDFOR
  \STATE \textbf{return} $y^{(t)}$
  \end{algorithmic}
  \end{algorithm}
